\section{Receiving the Dogmas without Reduction or Excess}

Dogma asks for faith because God reveals and the Church proposes, not because every believer has independently reconstructed the whole history. Theology nevertheless serves faith by showing coherence, answering objections, and distinguishing what is binding from what remains open. The result should be neither credulity nor embarrassed minimalism, but intelligent ecclesial assent ordered to worship and discipleship.

\subsection{The assent of faith}

The four Marian dogmas belong to truths to be believed with divine and Catholic faith. Their formal motive is the authority of God revealing, known through the Church's definitive proposal. That assent is not faith in Mary as a source of revelation and not an opinion that a historical argument happens to make more probable. It is faith in God concerning what he has done in her.

Historical and theological arguments have supporting roles. They show that the definitions are not arbitrary, disclose continuity, clarify terms, and remove apparent contradiction. They also expose the limits of bad arguments. A believer need not treat an apocryphal Dormition narrative, a disputed patristic quotation, an obsolete embryology, or a private apparition as the ground of faith. Indeed, attaching dogma to weak evidence can make faith appear to fall when the weak claim is corrected.

The exact proposition, not every theological explanation surrounding it, receives dogmatic assent. Catholics may discuss the best exegesis of a typological passage, the historical route of a feast, the authenticity of a sermon, Mary's death before the Assumption, or how best to describe her intercession, while retaining the defined objects. Legitimate plurality ends where a position contradicts what the Church has definitively proposed.

\subsection{Why definitions can be later than the apostolic age}

The date of definition is not the alleged date of invention. \emph{Homoousios}, the two-natures formula, and the canon's full recognition likewise required controversy and judgment. A later Marian definition claims that the truth belongs to the deposit through Scripture and Tradition, not that the apostles privately dictated the later juridical sentence.

At the same time, ``development'' cannot be used as a magic word. A credible development must preserve the Christological and soteriological center, account for earlier witnesses and silences, explain conceptual advances, and receive the Church's authoritative judgment. The Immaculate Conception's route through preservative redemption and the Assumption's route through Paschal hope, liturgy, and consensus are different developments. Neither is demonstrated by the bare fact that devotion grew.

\subsection{Frequent misunderstandings}

\begin{fourstable}{Claim}{What is true}{What is false}{Corrective locus}
``Mary is Mother of God'' & The person born from her according to humanity is God the Son. & She causes the divine nature, precedes God, or is a goddess. & Luke 1:35, 43; Ephesus; Chalcedon; ST III, q.35, a.4.\\
``Mary is ever-virgin'' & Virginal conception, inviolate birth, and lifelong virginity form one vocation. & Marriage or ordinary childbirth is impure; Christ's birth was unreal; speculative anatomy is dogma. & Matt 1; Luke 1--2; Lateran 649; CCC 496--507; ST III, q.28.\\
``Mary is immaculately conceived'' & Christ preserves her from original sin from her first instant. & Jesus rather than Mary is the subject; Mary has no need of Christ; Anne conceives virginally. & \emph{Ineffabilis Deus}; Rom 5; Eph 1; MPF 14.\\
``Mary is assumed'' & God glorifies her body and soul in anticipation of the common resurrection. & She ascends by her own power; Scripture narrates the event; the dogma defines whether she died. & 1 Cor 15; \emph{Munificentissimus Deus} 44; LG 59, 68.\\
``Mary is sinless'' & Singular grace preserves her from original and personal sin. & She is divine, omniscient, incapable of growth, or outside redemption. & Trent V and VI; LG 56; CCC 490--493.\\
``Mary intercedes'' & Her maternal prayer participates in the communion of saints under Christ. & The Trinity is reluctant, Mary changes a merciless Son, or grace must flow through a separate treasury. & John 2; LG 60--62; MPF 38--55.\\
``Mary is queen'' & She shares her Son's royal glory analogically and dependently. & She possesses equal sovereignty or natural omnipotence. & \emph{Ad caeli Reginam} 31, 34, 39; LG 59.\\
``Co-redemptrix'' & Advocates often intended subordinate cooperation. & The title is a fifth dogma or presently suitable constructive terminology. & MPF 17--22.\\
\end{fourstable}

\subsection{Veneration and worship}

Catholic tradition distinguishes adoration owed to the triune God from veneration given to saints. Mary's unique relation to Christ grounds a singular veneration, conventionally called \emph{hyperdulia}; it never becomes \emph{latria}. Vatican II says Marian devotion differs essentially from the adoration offered to the Word incarnate, Father, and Spirit, and rightly fosters that worship (LG 66). The distinction is not merely a private intention hidden beneath identical acts. Christian images, feasts, prayers, and preaching must make the dependence visible.

Calling Mary blessed fulfills Luke 1:48 when it magnifies the Mighty One's deed. Asking her prayer belongs to the communion of saints when it ends in trust in God. Marian practice is disordered when it eclipses Scripture and sacrament, assigns mechanical power to an object or formula, nourishes contempt for non-devotees, exploits fear, or replaces justice, medicine, safeguarding, and ordinary Christian responsibility.

\subsection{Ecumenical truthfulness}

The first two dogmas belong prominently to the common patristic inheritance, though communities differ over authority and implications. Eastern Orthodox liturgy and theology richly confess \emph{Theotokos}, \emph{Aeiparthenos}, and Dormition, while the Latin analysis of original sin and the juridical form of the 1854 and 1950 papal definitions create distinct ecumenical questions. Protestant communities differ among themselves; many retain the ancient Christological councils and virginal conception while rejecting later Marian definitions or invocation.

Catholic exposition should neither conceal these disagreements nor use Marian doctrine as a tribal test detached from Christ. It should state the Catholic claim exactly, identify common sources, acknowledge real differences in theological vocabulary and ecclesial authority, and avoid attributing caricatures to dialogue partners. MPF itself identifies ecumenical responsibility as one reason to prefer language that preserves the evident centrality of Christ.

Ecumenical concern cannot mean treating a dogma as optional. It means showing why Catholics believe it and how its true object excludes the excess an interlocutor may fear. \emph{Theotokos} protects Christ, perpetual virginity manifests divine initiative, immaculate preservation exalts Christ's merits, and Assumption confesses his resurrection victory. The deepest Catholic answer to the charge of ``too much Mary'' is not less exact doctrine but more exact Christology.

\subsection{Prayer shaped by the dogmas}

The dogmas can form four movements of Christian prayer:

\begin{enumerate}
  \item Before the Mother of God, confess Jesus Christ as one Lord, truly divine and truly human, and ask for faith that does not divide doctrine from life.
  \item Before the Ever-Virgin, honor God's initiative and ask for an undivided heart faithful to one's actual vocation, whether marriage, consecration, ordained ministry, or single discipleship.
  \item Before the Immaculate one, praise prevenient mercy and ask to receive grace as gift, without despair about sin or presumption about one's power.
  \item Before the Assumed one, hope for resurrection, honor the body, accompany the dying, bury the dead, and refuse every spirituality that treats creation as disposable.
\end{enumerate}

Mary's own scriptural speech remains the norm. She receives the word, magnifies the Lord, remembers covenant mercy, and tells servants to obey Christ (Luke 1:38, 46--55; John 2:5). Authentic dogmatic devotion does the same.

\subsection{Conclusion: the masterpiece belongs to the Artist}

The four dogmas make strong claims about a creature because they make stronger claims about God. The eternal Son can truly enter history without ceasing to be God. The Spirit can bring forth a new creation without violating created integrity. Christ's redemption can reach a person preveniently and perfect freedom rather than compete with it. The risen Lord can bring the whole human person into glory.

Mary's excellence neither fills a divine lack nor negotiates a compromise between justice and mercy. She is the Mother because the Word became flesh; ever-virgin because God's calling embraces the whole person; immaculate because the Savior saves perfectly; assumed because the risen Christ conquers death. She is first among the redeemed, never beside the Redeemer as his equal or supplement. To confess her rightly is to let her Magnificat govern theology: all generations call her blessed because the Mighty One has done great things, and holy is his name.
